\documentclass[10pt]{article}
\usepackage[margin=0.9in]{geometry}
\usepackage{amsmath,amssymb,amsthm,hyperref}
\newtheorem{theorem}{Theorem}
\title{The row-sum polynomials of Barry's almost Riordan array}
\author{Principia Math in Progress}
\date{13 September 2026}
\begin{document}
\maketitle

\paragraph{Task 1. Prove the row-sum formula.}
Conjecture 27 of Barry \cite{Barry} concerns the array produced by
\(P=\overline A\), where the bar deletes the first row and
\[
 A=\left(\frac{1+2z}{1-z^2};
          \frac{-1}{(1-z)(1+z)^2},\frac{z}{1+z}\right).
\]
In this almost Riordan notation the first column has generating
function \(a(z)=(1+2z)/(1-z^2)\), while column \(k\ge1\) has
generating function
\[
 z\frac{-1}{(1-z)(1+z)^2}\left(\frac{z}{1+z}\right)^{k-1}.
\]
Define the produced rows by \(u_0=(1,0,\ldots)\) and
\(u_{n+1}=u_nP\), and set \(P_n(x)=\sum_k(u_n)_kx^k\).
These operations are finite in each row because \(P\) is lower
Hessenberg.

\begin{theorem}
For every integer \(n\ge0\),
\begin{equation}\label{answer}
 P_n(x)=1+(1-x)\sum_{i=0}^{n-1}(1-x)^{\overline i},
 \qquad y^{\overline i}=\prod_{j=0}^{i-1}(y+j),
 \quad y^{\overline0}=1.
\end{equation}
\end{theorem}

\paragraph{Factorial convention.}
The rising factorial in \eqref{answer} is essential: it gives the
source's \(P_3=5-8x+5x^2-x^3\). The falling-factorial gloss added
in the MathDB transcription would give \(3-4x+3x^2-x^3\), which
already contradicts the source's displayed values. We prove the
full source formula, rather than counting that transcription
error as a counterexample.

\begin{proof}
Let \(R_m(x)=\sum_k A_{m,k}x^k\). The specified column generating
functions give, as a formal power series in \(z\),
\begin{align*}
 \sum_{m\ge0}R_m(x)z^m
 &=\frac{1+2z}{1-z^2}
   -\frac{zx}{(1-z)(1+z)(1+(1-x)z)}\\
 &=\frac{1+2(1-x)z}{(1-z)(1+(1-x)z)}\\
 &=\frac1{1-z}+
    \frac{(1-x)z}{(1-z)(1-(x-1)z)}.
\end{align*}
Consequently
\[
 R_{j+1}(x)=1+(1-x)\sum_{i=0}^{j}(x-1)^i.
\]
Since row \(j\) of \(P\) is row \(j+1\) of \(A\), right
multiplication by \(P\) acts on row polynomials by the linear
operator \(T\) defined by \(T(x^j)=R_{j+1}(x)\).
For every polynomial \(F\), summing the finite geometric series
therefore yields
\begin{equation}\label{operator}
 (TF)(x)=\frac{(3-2x)F(1)+(1-x)^2F(x-1)}{2-x}.
\end{equation}
The numerator vanishes at \(x=2\), so this is a polynomial
identity with no excluded value of \(x\). In particular,
\(P_{n+1}=TP_n\).

Write \(Q_n\) for the right side of \eqref{answer}.
We have \(Q_0=1=P_0\) and \(Q_n(1)=1\). Substitution into
\eqref{operator} gives
\begin{align*}
 TQ_n
 &=\frac{3-2x+(1-x)^2}{2-x}
      +(1-x)^2\sum_{i=0}^{n-1}(2-x)^{\overline i}\\
 &=2-x+(1-x)^2\sum_{i=0}^{n-1}(2-x)^{\overline i}\\
 &=1+(1-x)\left(1+
       \sum_{i=1}^{n}(1-x)^{\overline i}\right)=Q_{n+1}.
\end{align*}
The last step uses
\((1-x)^{\overline{i+1}}=(1-x)(2-x)^{\overline i}\).
Induction proves the formula for every \(n\).
\end{proof}

\paragraph{Task 2. Determine an explicit coefficient formula.}
Express each coefficient of \(P_n(x)\) through standard Stirling
numbers and obtain a combinatorial interpretation of its absolute
value.

\paragraph{Task 3. Analyze the zeros.}
Determine which \(P_n\) have only real zeros, and describe the
limiting zero distribution after an appropriate rescaling.

\paragraph{Task 4. Extend the production operator.}
Identify parameter families of almost Riordan arrays for which
the induced operator has an analogous evaluation-and-shift
formula and the produced polynomials admit finite factorial sums.

\paragraph{Task 5. Derive asymptotics.}
Obtain asymptotic formulas for \(P_n(x)\) at fixed real or complex
\(x\), including the exceptional values where the factorial sum
terminates or cancels.

\begin{thebibliography}{1}
\bibitem{Barry} P.~Barry, \emph{On the partial sums of Riordan
arrays}, \href{https://arxiv.org/abs/2108.13537}
{arXiv:2108.13537v1} (2021). The almost Riordan array is on
pp.~19--20, the produced array and its initial polynomials on
p.~21, and Conjecture 27 on p.~22 of the inspected 23-page PDF.
\end{thebibliography}
\end{document}
